% ============================================================
% ENTITY LOCATION REFERENCES
% Viability Philosophy System
% Book I: Ontology of Viability
% ============================================================

\cleardoublepage

\section*{Entity Location References}

\phantomsection
\addcontentsline{toc}{section}{Entity Location References}

% ------------------------------------------------------------
% 1. PURPOSE AND SCOPE
% ------------------------------------------------------------

\subsection*{1. Purpose and Scope}

Entity location references establish the connection
between logical entities and the structural locations
at which those entities occur within the text of the
Viability Philosophy System.

A logical entity identifier establishes the identity
of an entity independently of its textual location.
A structural identifier establishes the location of
a textual element within the hierarchy of a book.

Entity location references connect these two forms
of identification.

Their purpose is to make it possible to determine
where a particular logical entity is introduced,
defined, examined, qualified, applied, or otherwise
discussed within the manuscript.

An entity may occur in several textual locations.
These locations may serve different functions in
the development of the discussion, but they do not
necessarily constitute different logical entities.

The entity location referencing system therefore
preserves the distinction between the identity of
an entity and the plurality of its textual occurrences.

It supports:

\begin{itemize}
    \item locating the textual occurrences of a
    logical entity;
    \item distinguishing an entity's initial
    presentation from its subsequent discussions;
    \item tracing the development of an entity
    across different parts of the book;
    \item identifying the passages relevant to
    the interpretation or application of an entity;
    \item coordinating logical entity inventories
    with the structural organization of the manuscript;
    \item maintaining accurate location records
    during revision.
\end{itemize}

% ------------------------------------------------------------
% 2. THE RELATION BETWEEN ENTITY IDENTITY
% AND TEXTUAL LOCATION
% ------------------------------------------------------------

\subsection*{2. The Relation Between Entity Identity
and Textual Location}

A logical entity and its textual location
are different kinds of objects of identification.

The logical entity identifier designates
the entity within the logical organization
of the system.

The structural identifier designates
a particular location within the textual
organization of a book.

For example:

\begin{center}
\begin{tabular}{ll}
    \toprule
    \textbf{Identifier} & \textbf{Function} \\
    \midrule
    \texttt{AX-001}
        & Logical entity identity \\
    \texttt{VS.B01.P01.C01.S01.SS01.P001}
        & Structural location \\
    \bottomrule
\end{tabular}
\end{center}

The first identifier remains associated
with the same logical entity across
its textual occurrences.

The second identifier identifies
one particular location in which
that entity may be discussed.

Consequently, one logical entity may
be associated with multiple structural
locations.

Conversely, a single structural location
may discuss, introduce, or relate
to more than one logical entity.

The relation between entities and
locations is therefore not necessarily
one-to-one.

% ------------------------------------------------------------
% 3. BASIC PRINCIPLE OF ENTITY LOCATION
% ------------------------------------------------------------

\subsection*{3. Basic Principle of Entity Location}

Every recorded entity location must
identify both sides of the relation:

\begin{enumerate}
    \item the logical entity to which
    the location belongs;
    \item the structural location at
    which the entity occurs.
\end{enumerate}

The logical entity must be identified
through its assigned logical entity
identifier.

The textual location must be identified
through its corresponding structural
identifier.

A location record must not replace
the logical entity identifier with
a textual name alone when an established
identifier is available.

Likewise, a logical entity identifier
must not be treated as sufficient
information for locating a particular
textual occurrence.

The two identifiers must be used
together whenever the purpose is
to establish the location of an
entity within the manuscript.

% ------------------------------------------------------------
% 4. MULTIPLE LOCATIONS OF A SINGLE ENTITY
% ------------------------------------------------------------

\subsection*{4. Multiple Locations of a Single Entity}

A logical entity may be discussed
in several locations throughout
the book.

Its treatment may extend across
different chapters, sections,
subsections, and paragraphs.

The existence of multiple locations
does not, by itself, imply the
existence of multiple entities.

For example, an entity identified
as \texttt{AX-001} may be introduced
in one paragraph, examined in
another, and used in a later
argument.

Each occurrence has its own
structural location, while the
logical entity retains its original
identifier.

The location referencing system
must preserve this distinction.

Where an entity occurs in several
locations, the relevant locations
should be recorded without
duplicating the entity's logical
identity.

% ------------------------------------------------------------
% 5. FUNCTIONS OF TEXTUAL OCCURRENCES
% ------------------------------------------------------------

\subsection*{5. Functions of Textual Occurrences}

Different textual occurrences
of the same logical entity may
serve different functions in
the development of the book.

The function of an occurrence
should be distinguished from
the identity of the entity itself.

Depending on the actual content
of a passage, an occurrence may
serve one or more of the following
descriptive functions.

\subsubsection*{5.1. Introduction}

A passage may introduce an entity
into the discussion.

The corresponding location helps
the reader identify where the
entity first becomes relevant
within the textual development.

\subsubsection*{5.2. Definition or Characterization}

A passage may provide a formal
definition or a substantive
characterization of an entity.

The location should be recorded
when identifying the passage
is important for subsequent
interpretation or reference.

An introductory mention must
not automatically be treated
as a formal definition.

\subsubsection*{5.3. Explanation or Analysis}

An entity may be examined through
explanatory, analytical, or
interpretive discussion.

Such passages may clarify its
features, scope, conditions,
or relations to other entities.

Their locations should remain
associated with the same logical
entity when the identity of
that entity is preserved.

\subsubsection*{5.4. Qualification or Revision}

A passage may qualify, restrict,
or revise an earlier account
of an entity.

The corresponding location
should be recorded when the
qualification or revision
is relevant to understanding
the entity's development.

A revision of the account must
be distinguished from a change
in the identity of the entity.

\subsubsection*{5.5. Application or Use}

An entity may be used in an
argument, explanation, or
discussion without being
redefined at that location.

Such an occurrence may be
recorded to make the use
of the entity traceable.

\subsubsection*{5.6. Comparison or Relation}

A passage may discuss an entity
in comparison with, or in relation
to, another entity.

The location may be associated
with each relevant entity,
provided that the passage
actually addresses them.

The association of several
entities with one location
does not imply that those
entities are identical or
that a particular logical
relation has been established.

% ------------------------------------------------------------
% 6. PRIMARY AND ADDITIONAL LOCATIONS
% ------------------------------------------------------------

\subsection*{6. Primary and Additional Locations}

Where an entity is discussed
in multiple places, it may
be useful to distinguish
the location most central
to its formal presentation
from other locations in
which it appears.

For practical referencing
purposes, a location record
may distinguish between:

\begin{itemize}
    \item a primary location,
    where the entity receives
    its principal presentation
    or formal treatment;

    \item additional locations,
    where the entity is further
    explained, examined,
    qualified, applied, or
    related to other entities.
\end{itemize}

The designation of a primary
location is a navigational
and organizational convention.

It does not establish that
the entity exists only at
that location, nor does it
imply that every other
occurrence is conceptually
secondary.

The primary location should
be selected according to
the actual organization and
function of the manuscript.

Where no single location
can appropriately serve
as the primary location,
the inventory should not
force an artificial choice.

% ------------------------------------------------------------
% 7. PRECISION OF LOCATION REFERENCES
% ------------------------------------------------------------

\subsection*{7. Precision of Location References}

An entity location reference
should identify the narrowest
structural unit necessary
to locate the relevant
discussion accurately.

The appropriate level of
specificity depends on
the function of the passage.

A chapter-level reference
may be sufficient when
the entity is treated
throughout a chapter.

A section- or subsection-level
reference may be preferable
when the relevant discussion
is more concentrated.

A paragraph-level reference
should be used when a specific
passage must be located
precisely and an appropriate
paragraph identifier exists.

The referencing system must
not imply a degree of precision
that the structural organization
of the manuscript does not
support.

A location reference must
correspond to an actual
structural element.

% ------------------------------------------------------------
% 8. ENTITY LOCATION RECORDS
% ------------------------------------------------------------

\subsection*{8. Entity Location Records}

The relationship between
logical entities and textual
locations should be maintained
in a consistent and traceable
record.

At a minimum, such a record
must make it possible to
identify:

\begin{itemize}
    \item the logical entity
    identifier;
    \item the structural
    identifier of the
    relevant location;
    \item the descriptive
    function of the occurrence,
    where this information
    is necessary;
    \item whether the location
    is designated as primary
    or additional, where
    such a distinction is used.
\end{itemize}

A location record may also
contain a brief description
of the relevant passage
when this improves the
accuracy of retrieval.

Descriptions must remain
informational. They must
not substitute for the
formal identifiers.

The record should be maintained
in coordination with the
logical entity inventory
and the structural organization
of the manuscript.

% ------------------------------------------------------------
% 9. RELATION TO INTERNAL
% CROSS-REFERENCES
% ------------------------------------------------------------

\subsection*{9. Relation to Internal
Cross-References}

Entity location references
and internal cross-references
are closely related, but
they are not identical.

An internal cross-reference
directs the reader from
one textual location or
logical entity to another.

An entity location reference
records the association
between a logical entity
and one or more locations
in the manuscript.

The former primarily serves
navigation within the text.

The latter establishes
traceability between the
logical inventory and the
textual occurrences of
its entities.

A cross-reference may use
an entity location record
to direct the reader to
a relevant passage.

However, the existence
of a location record does
not, by itself, determine
which cross-reference
should appear in a given
passage.

The selection of a cross-reference
must remain appropriate
to the argument and context.

% ------------------------------------------------------------
% 10. RELATION TO LOGICAL
% DEPENDENCIES
% ------------------------------------------------------------

\subsection*{10. Relation to Logical
Dependencies}

Entity location references
must be distinguished from
logical dependencies.

An entity location reference
establishes where an entity
is discussed.

A logical dependency specifies
a relation between logical
entities within the system.

The fact that two entities
occur in the same paragraph,
section, or chapter does
not establish that one
depends on the other.

Likewise, the fact that
one entity is discussed
after another does not
establish a relation of
derivation, entailment,
or logical priority.

When a dependency between
entities is intended,
it must be specified through
the logical dependency
system.

Location records may help
the reader find the passages
in which a dependency is
discussed, but they do not
constitute the dependency
itself.

% ------------------------------------------------------------
% 11. RELATION TO STRUCTURAL
% AND LOGICAL IDENTIFIERS
% ------------------------------------------------------------

\subsection*{11. Relation to Structural
and Logical Identifiers}

Entity location references
depend on the consistent
use of two identification
systems.

Structural identifiers
locate textual elements.

Logical entity identifiers
identify entities.

The entity location system
connects them without
collapsing their distinct
functions.

The structural identifier
must reflect the current
location of the passage.

The logical entity identifier
must reflect the identity
of the entity.

A change in one does not
automatically require a
change in the other.

For example, moving a passage
to another chapter may
change its structural
identifier while leaving
the logical identity of
the entity discussed
in that passage unchanged.

Conversely, a conceptual
replacement may require
a new logical entity
identifier even if the
new entity is introduced
at the same structural
location.

% ------------------------------------------------------------
% 12. MAINTENANCE DURING
% MANUSCRIPT REVISION
% ------------------------------------------------------------

\subsection*{12. Maintenance During
Manuscript Revision}

Entity location records
must be reviewed whenever
the manuscript is revised
in a way that affects
the location or treatment
of a logical entity.

The following cases require
particular attention.

\subsubsection*{12.1. Relocation of a Passage}

When a passage is moved,
its structural identifier
may change.

All entity location records
associated with that passage
must be reviewed and updated
accordingly.

The logical entity identifiers
must remain unchanged when
the identities of the entities
have been preserved.

\subsubsection*{12.2. Addition of a New Occurrence}

When a new passage introduces,
examines, or uses an existing
entity, the new location
should be associated with
the existing logical entity
identifier.

A new textual occurrence
does not automatically
require a new logical
entity identifier.

\subsubsection*{12.3. Removal of a Passage}

When a passage is removed,
its location record must
be reviewed.

If the passage was the
only recorded location
of an entity, the entity
inventory must also be
examined to determine
whether the entity remains
part of the system and
whether it is presented
elsewhere.

\subsubsection*{12.4. Revision of an Entity}

When the conceptual treatment
of an entity changes, its
location records should
be reviewed to ensure
that they continue to
represent the actual
content of the relevant
passages.

A change in exposition
must not automatically
be interpreted as a change
in entity identity.

\subsubsection*{12.5. Change of Primary Location}

If a different passage
becomes the principal
location for an entity's
formal presentation,
the primary location
designation should be
updated.

The previous location
may remain recorded
as an additional location
if it continues to be
relevant.

% ------------------------------------------------------------
% 13. EXTERNAL SOURCES
% ------------------------------------------------------------

\subsection*{13. External Sources}

Entity location references
concern the locations of
internal logical entities
within the Viability
Philosophy System.

They do not assign internal
logical entity identifiers
to external books, articles,
authors, documents, or
other bibliographic sources.

An external source may
be cited or discussed
at a particular structural
location.

That location belongs
to the internal structure
of the book, while the
external source retains
its own bibliographic
identity.

If an internal logical
entity is discussed in
relation to an external
source, the entity and
the source must remain
separately identifiable.

The relevant external
reference conventions
are established in the
documents governing
external source references
and bibliographic organization.

% ------------------------------------------------------------
% 14. SUMMARY
% ------------------------------------------------------------

\subsection*{14. Summary}

Entity location references
connect logical entity
identities with their
textual occurrences.

They operate through
the coordination of:

\begin{enumerate}
    \item logical entity
    identifiers, which
    identify the entities;

    \item structural
    identifiers, which
    identify the locations
    in which those entities
    occur.
\end{enumerate}

A single logical entity
may be associated with
multiple structural
locations, and a single
structural location may
be associated with
multiple logical entities.

The system must preserve
the distinction between
entity identity, textual
occurrence, and logical
dependency.

Its records must remain
accurate and be updated
whenever manuscript
revisions affect the
locations or treatment
of the entities.

Entity location references
therefore provide the
traceability layer connecting
the logical organization
of the Viability Philosophy
System with the textual
architecture of its books.

\cleardoublepage

% ============================================================
% END OF ENTITY LOCATION REFERENCES
% ============================================================